\documentclass[11pt,a4paper]{article}
\usepackage[utf8]{inputenc}
\usepackage[T1]{fontenc}
\usepackage{amsmath, amssymb}
\usepackage{amsthm}
\usepackage[margin=2.5cm]{geometry}
\usepackage{hyperref}
\hypersetup{pdftitle={Energy Decay in a Damped Harmonic Oscillator}, pdfauthor={Lena Ortiz}}

\newtheorem{theorem}{Theorem}

\title{Energy Decay in a Damped Harmonic Oscillator}
\author{Lena Ortiz}
\date{\today}

\begin{document}
\maketitle

\begin{equation}
  m\ddot{\{x\}}+c\dot{\{x\}}+kx=0,
\end{equation}
where m \textgreater{} 0 is the mass, c ≥ 0 the damping coefficient and k \textgreater{} 0 the spring constant. Dividing by m and writing omega\_0 = sqrt(k/m) and gamma = c/(2m) gives the standard form
\begin{equation}
  \ddot{\{x\}}+2\gamma\dot{\{x\}}+\omega_{0}^{2}x=0.
\end{equation}
2 Solution\\

Substituting x = exp(lambda t) into equation (2) gives the characteristic equation lambda\textasciicircum{}2 + 2 gamma lambda + omega\_0\textasciicircum{}2 = 0, whose roots are
\begin{equation}
  \lambda=-\gamma\pm\sqrt{\gamma^{2}-}\omega_{0}^{2}.
\end{equation}
The behaviour depends on the sign of the discriminant:
\\

\{underdamped gamma \textless{} 0
critically damped gamma = 0
overdamped gamma \textgreater{} 0\}
\\

In the underdamped case let omega\_d = sqrt\{omega\_0\textasciicircum{}2 - gamma\textasciicircum{}2\}. With x(0) = x\_0 and dot x (0) = 0 the solution is
\begin{equation}
  x(t)=x_{0}\exp\{-\gamma\}t(\cos\{\omega_{d}t\})+\frac{\gamma}{\omega_{d}}\sin\{\omega_{d}t\}
\end{equation}
Equivalently, with y = (x, dot x) the second-order equation becomes the first-order system
\begin{equation}
  \dot{\{y\}}=\begin{pmatrix}
0;1;-\omega_{0}^{2};-2\gamma &  \\
y & 
\end{pmatrix}
\end{equation}

\begin{theorem}
  The energy E(t) = frac\{1\}\{2\} m dot\{x\}\textasciicircum{}2 + frac\{1\}\{2\} k x\textasciicircum{}2 never increases. More precisely, dE/dt = -c dot\{x\}\textasciicircum{}2 \textless{}= 0.
\end{theorem}

\begin{proof}
  Differentiate E and use (1):\\
\\
dE/dt = m dot\{x\} ddot\{x\} + k x dot\{x\}\\
 = dot\{x\}(m ddot\{x\} + k x)\\
 = -c dot\{x\}\textasciicircum{}2.
\end{proof}

3 Numerical check\\

We integrated (2) with omega\_0 = 2, gamma = 0.2 and x\_0 = 1 on t in [0, 10], using the explicit Euler method and the classical fourth-order Runge-Kutta method (RK4) [2]. For each step size h:
\\

1. integrated from t = 0 to t = 10 with 10/h steps;
2. compared every step with the exact solution (4);
3. recorded the largest absolute difference.
\\

References
\\

[1] H. Goldstein, C. Poole and J. Safko, Classical Mechanics, 3rd ed., Addison-Wesley, 2002.
\\

[2] E. Hairer, S. P. Norsett and G. Wanner, Solving Ordinary Differential Equations I, 2nd ed., Springer, 1993.

\end{document}
